\section{Объектно-ориентированное проектирование и диаграммы}
\subsection{Объектный анализ}
Предметные сущности --- отдел, сотрудник, помещение и мероприятие. Данные редактируемого плана выделены в \code{EventData}. \code{CalendarEvent} добавляет идентификатор, статус, версии и снимки сведений об участниках. Изменение данных отделено от отображения: графический объект сотрудника связан с моделью по ID, но не управляет её временем.

\code{CalendarService} хранит календарь.\par
\code{SimulationEngine} задаёт ход эксперимента. Планирование уведомлений, подбор переносов, построение отчёта и сохранение представлены отдельными классами. Композиция используется для объектов одного прогона. Наследование введено в канале уведомлений: абстрактный \code{NotificationChannel} реализуется классом \code{SimulatedEmailChannel}. Классы интерфейса наследуются от \code{QWidget}, \code{QDialog} или \code{QObject}.

На диаграммах сплошная линия с ромбом показывает композицию, пунктирная стрелка --- использование, пустой треугольник --- наследование. В диаграммах перечислены основные поля и операции, а не весь интерфейс Qt и вспомогательные типы.

\begin{figure}[H]\centering
\resizebox{\linewidth}{!}{\begin{tikzpicture}
\node[uml] (catalog) at (0,0) {\textbf{CompanyCatalog}\par\hrule\smallskip departments: QMap\par employees: QMap\par rooms: QMap};
\node[uml] (department) at (-5,-3) {\textbf{Department}\par\hrule\smallskip id: int\par name: QString\par leaderId: int};
\node[uml] (employee) at (0,-3) {\textbf{Employee}\par\hrule\smallskip id: int\par name: QString\par departmentId: int\par publicId: QString};
\node[uml] (room) at (5,-3) {\textbf{Room}\par\hrule\smallskip id: int\par name: QString\par departmentId: optional<int>};
\draw[composition] (catalog.south west)--(department.north) node[midway,left] {5..9};
\draw[composition] (catalog.south)--(employee.north) node[midway,right] {1..*};
\draw[composition] (catalog.south east)--(room.north) node[midway,right] {K+1};
\node[uml] (event) at (-2.5,-6.7) {\textbf{CalendarEvent}\par\hrule\smallskip id, version, scheduleVersion\par status: EventStatus\par data: EventData\par participantCodes, organizerCode};
\node[uml] (data) at (2.5,-6.7) {\textbf{EventData}\par\hrule\smallskip title, type, start, end\par roomId: optional<int>\par participantIds: QVector<int>\par organizerId, importance};
\draw[composition] (event.east)--(data.west) node[midway,above] {1};
\draw[dep] (data.north)--(employee.south) node[pos=.35,below,sloped,fill=white,inner sep=1pt,font=\scriptsize] {участники};
\draw[dep] (data.north east)--(room.south) node[midway,right] {ID};
\end{tikzpicture}}

\caption{Диаграмма классов предметной области. Ссылки на сотрудников и помещения хранятся как ID.}
\end{figure}
\begin{figure}[H]\centering
\resizebox{\linewidth}{!}{\begin{tikzpicture}
\node[uml] (engine) at (0,0) {\textbf{SimulationEngine}\par\hrule\smallskip setup, random, requests, log\par step(), reset(), trigger()\par createEvent(), updateEvent()\par acceptSuggestion()};
\node[uml] (calendar) at (-5,-4) {\textbf{CalendarService}\par\hrule\smallskip catalog, active, archive, history\par checkEvent(), createEvent()\par updateEvent(), calendar()\par advanceTo()};
\node[uml] (reminders) at (0,-4) {\textbf{ReminderService}\par\hrule\smallskip pending: QMap\par observe(records, rules)\par deliver(calendar, channel)\par clear()};
\node[uml] (email) at (5,-4) {\textbf{SimulatedEmailChannel}\par\hrule\smallskip messages: QVector\par send(EmailMessage)\par messages(), clear()};
\draw[composition] (engine.south west)--(calendar.north);
\draw[composition] (engine.south)--(reminders.north);
\draw[composition] (engine.south east)--(email.north);
\node[uml] (channel) at (5,-7.4) {\textit{NotificationChannel}\par\hrule\smallskip virtual send(message) = 0\par virtual destructor};
\draw[-{Triangle[open,length=3mm]},thick] (email.south)--(channel.north);
\draw[dep] (reminders.south)--(channel.west) node[midway,below,sloped] {использует};
\node[uml] (resolver) at (-5,-7.4) {\textbf{ConflictResolver}\par\hrule\smallskip static alternatives(calendar,\par setup, requested, actor, id)};
\draw[dep] (engine.west)--++(-3.1,0)--++(-4.4,0)|-(resolver.west);
\draw[dep] (resolver.north)--(calendar.south);
\node[uml] (report) at (-5,-10.6) {\textbf{SimulationReport}\par\hrule\smallskip create(), metrics(), csv()\par summary()};
\node[uml] (store) at (0,-10.6) {\textbf{SimulationStore}\par\hrule\smallskip snapshot(), restore()\par save(), load()};
\node[uml] (runner) at (5,-10.6) {\textbf{ExperimentRunner}\par\hrule\smallskip current: unique\_ptr<Engine>\par comparison(), advance()\par report(), csv()};
\end{tikzpicture}}
\par\smallskip{\footnotesize SimulationReport читает движок; SimulationStore воспроизводит его команды; ExperimentRunner создаёт независимый движок каждого прогона. Эти зависимости вынесены в подпись, чтобы не перегружать линии диаграммы.}

\caption{Основные классы ядра и иерархия канала уведомлений.}
\end{figure}
\begin{figure}[H]\centering
\resizebox{\linewidth}{!}{\begin{tikzpicture}
\node[uml] (qwidget) at (-5,0) {\textit{QWidget}};
\node[uml] (qdialog) at (0,0) {\textit{QDialog}};
\node[uml] (qobject) at (5,0) {\textit{QObject}};
\node[uml] (widget) at (-5,-2.4) {\textbf{Widget}\par\hrule\smallskip параметры, роль\par открытие нового прогона\par загрузка и серия опытов};
\node[uml] (window) at (0,-2.4) {\textbf{SimulationWindow}\par\hrule\smallskip engine: SimulationEngine\par timer: QTimer\par календарь и ручные действия};
\node[uml] (animator) at (5,-2.4) {\textbf{OfficeAnimator}\par\hrule\smallskip scene: QGraphicsScene\par update(), pause(), resume()\par finish(), retime()};
\draw[-{Triangle[open,length=3mm]}] (widget.north)--(qwidget.south);
\draw[-{Triangle[open,length=3mm]}] (window.north west)--(qwidget.south east);
\draw[-{Triangle[open,length=3mm]}] (animator.north)--(qobject.south);
\node[uml] (tools) at (-5,-5.5) {\textbf{RunTools : QWidget}\par\hrule\smallskip заявки и выбор переноса\par письма и экспорт};
\node[uml] (dialog) at (0,-5.5) {\textbf{ExperimentDialog}\par\hrule\smallskip runner: ExperimentRunner\par прогресс, остановка, экспорт};
\node[uml] (engine) at (5,-5.5) {\textbf{SimulationEngine}\par\hrule\smallskip ядро, без Qt Widgets};
\draw[dep] (widget.east)--(window.west) node[midway,above] {создаёт};
\draw[composition] (window.south west)--(tools.north east);
\draw[composition] (window.south east)--(engine.north west);
\draw[dep] (window.east)--(animator.west);
\draw[-{Triangle[open,length=3mm]}] (dialog.north east)--++(0.8,0.4)--++(0,4.1)--(qdialog.east);
\end{tikzpicture}}

\caption{Классы представления и их связь с моделью. Владение элементами интерфейса определяется родителями Qt.}
\end{figure}

\subsection{Диаграмма объектов}
Рисунок~\ref{fig:objects} показывает иллюстративный снимок до первого шага: начало 12 октября 2026 года, шесть отделов, по восемь сотрудников, включён регулярный календарь. Объекты ниже являются экземплярами классов проекта. Имена \code{engine}, \code{calendar} и \code{email} обозначают роли объектов; это не отдельные глобальные переменные программы. Из каталога показан только один отдел и один сотрудник.
\begin{figure}[H]\centering
\resizebox{\linewidth}{!}{\begin{tikzpicture}
\node[uml] (engine) at (0,0) {\underline{\shortstack[l]{engine :\\SimulationEngine}}\par\hrule\smallskip seed = 42\par modelVersion = calendar-v7\par шаг = 30 минут};
\node[uml] (calendar) at (-5,-3.2) {\underline{\shortstack[l]{calendar :\\CalendarService}}\par\hrule\smallskip время = 12.10.2026 00:00\par active = исходные планы\par archive = пусто};
\node[uml] (reminders) at (0,-3.2) {\underline{\shortstack[l]{reminders :\\ReminderService}}\par\hrule\smallskip pending = сроки напоминаний\par для начальных планов};
\node[uml] (email) at (5,-3.2) {\underline{\shortstack[l]{email :\\SimulatedEmailChannel}}\par\hrule\smallskip messages = пусто\par отправки до первого шага нет};
\draw[composition] (engine.south west)--(calendar.north);
\draw[composition] (engine.south)--(reminders.north);
\draw[composition] (engine.south east)--(email.north);
\node[uml] (department) at (-5,-6.5) {\underline{d1 : Department}\par\hrule\smallskip id = 1\par name = «Отдел 1»\par leaderId = 1};
\node[uml] (employee) at (0,-6.5) {\underline{e1 : Employee}\par\hrule\smallskip id = 1\par departmentId = 1\par publicId = «1\_01»};
\node[uml] (event) at (5,-6.5) {\underline{p1 : CalendarEvent}\par\hrule\smallskip type = Planning\par start = 12.10.2026 09:00\par end = 09:30; roomId = 1\par status = Planned};
\draw[dep] (calendar.south)--(department.north) node[midway,left] {каталог};
\draw[dep] (department.south) -- ++(0,-0.75) -| (employee.south);
\node[font=\footnotesize,fill=white] at (-2.5,-8.4) {руководитель};
\draw[dep] (event.south) -- ++(0,-0.75) -| (employee.south);
\node[font=\footnotesize,fill=white] at (2.5,-8.4) {участник};
\draw[dep] (calendar.south east) -- ++(0,-0.4) -| (event.north);
\end{tikzpicture}}

\caption{Диаграмма объектов одного прогона (частичный снимок).}\label{fig:objects}
\end{figure}

\subsection{Правила модельного времени и актуализации}
Ресурсы заняты на полуоткрытом интервале $[s_i,e_i)$. События $i$ и $j$ пересекаются, если
\[
 s_i<e_j\quad\land\quad s_j<e_i.
\]
При пересечении общие участники или одинаковое непустое помещение создают конфликт. Касание границ, например 10:00--11:00 и 11:00--12:00, допустимо. Назначение в прошлое запрещено; произвольного пятичасового ожидания нет.

Состояния мероприятия: \code{Planned} $\rightarrow$ \code{InProgress} $\rightarrow$ \code{Completed}. Из \code{Planned} также возможна сервисная отмена с состоянием \code{Cancelled}. Завершённые и отменённые записи находятся в архиве.

Один шаг 30/60 минут последовательно обрабатывает каждую внутреннюю минуту:
\begin{enumerate}
\item Завершение событий и архивирование; старты на границе пока отложены.
\item В рабочее время --- создание периодических и случайных заявок, попытки изменения расписания.
\item Отправка наступивших имитируемых напоминаний по актуальному календарю.
\item Запуск мероприятий на текущей границе и публикация истории.
\end{enumerate}
На границе периода обрабатываются завершения, но новые заявки, уведомления и старты уже не создаются. Событие, начавшееся до конца периода и продолжающееся после него, остаётся \code{InProgress}; будущий план --- \code{Planned}. Они доступны в итоговом отчёте.

\subsection{Подбор разрешений конфликта}
Поиск сохраняет участников и длительность заявки и перебирает допустимые время и помещения. Используются рабочие дни в пределах заданного горизонта, начала на сетке 30 минут, исходное помещение, затем остальные помещения по ID. Каждый кандидат повторно проверяется календарным сервисом. Более раннее начало, чем исходно запрошенное, не предлагается. Поиск не переносит чужие мероприятия автоматически.

\code{Reject} фиксирует отказ. \code{Suggest} сохраняет несколько предложений до выбора пользователя. \code{AutoReschedule} применяет первый допустимый вариант немедленно. Если подходящих вариантов нет, заявка отклоняется. До принятия ожидающая заявка не резервирует ресурсы, поэтому вариант может устареть и требует повторной проверки.
